%-------------------------
% Resume in LaTeX — Karthikheyaa Kurra (research-oriented)
% Based on Jake Gutierrez's template (jakegut/resume), MIT license
%------------------------

\documentclass[letterpaper,11pt]{article}

\usepackage{latexsym}
\usepackage[empty]{fullpage}
\usepackage{titlesec}
\usepackage{marvosym}
\usepackage[usenames,dvipsnames]{color}
\usepackage{verbatim}
\usepackage{enumitem}
\usepackage[hidelinks]{hyperref}
\usepackage{fancyhdr}
\usepackage[english]{babel}
\usepackage{tabularx}
\input{glyphtounicode}

\pagestyle{fancy}
\fancyhf{}
\fancyfoot{}
\renewcommand{\headrulewidth}{0pt}
\renewcommand{\footrulewidth}{0pt}

% Adjust margins
\addtolength{\oddsidemargin}{-0.5in}
\addtolength{\evensidemargin}{-0.5in}
\addtolength{\textwidth}{1in}
\addtolength{\topmargin}{-.5in}
\addtolength{\textheight}{1.2in}

\urlstyle{same}

\raggedbottom
\raggedright
\setlength{\tabcolsep}{0in}

% Sections formatting
\titleformat{\section}{
  \vspace{-6pt}\scshape\raggedright\large
}{}{0em}{}[\color{black}\titlerule \vspace{-5pt}]

% Ensure that generated pdf is machine readable/ATS parsable
\pdfgentounicode=1

%-------------------------
% Custom commands
\newcommand{\resumeItem}[1]{
  \item\small{
    {#1 \vspace{-2pt}}
  }
}

\newcommand{\resumeSubheading}[4]{
  \vspace{-2pt}\item
    \begin{tabular*}{0.97\textwidth}[t]{l@{\extracolsep{\fill}}r}
      \textbf{#1} & #2 \\
      \textit{\small#3} & \textit{\small #4} \\
    \end{tabular*}\vspace{-7pt}
}

\newcommand{\resumeSubItem}[1]{\resumeItem{#1}\vspace{-4pt}}

\renewcommand\labelitemii{$\vcenter{\hbox{\tiny$\bullet$}}$}

\newcommand{\resumeSubHeadingListStart}{\begin{itemize}[leftmargin=0.15in, label={}]}
\newcommand{\resumeSubHeadingListEnd}{\end{itemize}}
\newcommand{\resumeItemListStart}{\begin{itemize}}
\newcommand{\resumeItemListEnd}{\end{itemize}\vspace{-7pt}}

%-------------------------------------------
%%%%%%  RESUME STARTS HERE  %%%%%%%%%%%%%%%%%%%%%%%%%%%%

\begin{document}

%----------HEADING----------
\begin{center}
    \textbf{\Huge \scshape Karthikheyaa Kurra} \\ \vspace{1pt}
    \small 267-299-9245 $|$ \href{mailto:karth@udel.edu}{karth@udel.edu} $|$
    \href{https://github.com/Khey17}{\underline{GitHub}} $|$
    \href{https://khey17.github.io/portfolio/}{\underline{Portfolio}} $|$
    \href{https://linkedin.com/in/karthikheyaa}{\underline{LinkedIn}}
\end{center}

%-----------EDUCATION-----------
\section{Education}
  \resumeSubHeadingListStart
    \resumeSubheading
      {University of Delaware}{Newark, DE}
      {M.S.\ Data Science, GPA 3.80/4.00}{Aug 2024 -- May 2026}
      \resumeItemListStart
        \resumeItem{\textbf{Coursework:} Probability Theory \& Applications, Mathematical Techniques for Data Science (Optimization), Applied Multivariate Statistics, Machine Learning, Algorithm Design \& Analysis, Research Software Engineering}
      \resumeItemListEnd
    \resumeSubheading
      {Sathyabama Institute of Science \& Technology}{Chennai, India}
      {B.E.\ Computer Science \& Engineering}{Oct 2020 -- May 2024}
  \resumeSubHeadingListEnd

%-----------EXPERIENCE-----------
\section{Experience}
  \resumeSubHeadingListStart

    \resumeSubheading
      {Research Associate, Scientific Machine Learning}{Jun 2026 -- Present}
      {Department of Mathematical Sciences, University of Delaware \, (Advisor: Dr.\ Ke Chen)}{Newark, DE}
      \resumeItemListStart
        \resumeItem{Writing a paper on hybrid PDE solvers that combine classical methods and neural networks to solve a nonlinear elliptic PDE in a non-iterative manner through learned surrogates and by leveraging hierarchical structure.}
        \resumeItem{Own end-to-end research infrastructure: high-precision ground-truth generation via Chebyshev spectral discretizations and Newton iteration, automated dataset-validation suites, and reproducible PyTorch training and evaluation pipelines.}
        \resumeItem{Drove architecture and loss design through controlled ablation studies covering network normalization and physics-informed residual losses with adaptive loss balancing, alongside systematic wall-clock profiling of classical versus learned inference.}
        \resumeItem{Maintain a regression-locked, unit-tested codebase with automated experiment reports that set weekly research direction with the PI.}
      \resumeItemListEnd

    \resumeSubheading
      {Research Assistant, Scientific Machine Learning}{Jun 2025 -- Jun 2026}
      {Department of Mathematical Sciences, University of Delaware \, (Advisor: Dr.\ Ke Chen)}{Newark, DE}
      \resumeItemListStart
        \resumeItem{Joined Dr.\ Ke Chen's group to build neural surrogates for nonlinear PDE solvers, replacing per-instance Newton iteration with amortized single-pass inference on new problem instances.}
        \resumeItem{Established spectral ground-truth generation, dataset validation, and reproducible PyTorch experiment pipelines that underpin the project's training and evaluation workflow.}
      \resumeItemListEnd

    \resumeSubheading
      {Research Intern, Computer Vision}{Mar 2023 -- Jun 2023}
      {Indian Institute of Technology, Hyderabad}{Hyderabad, India}
      \resumeItemListStart
        \resumeItem{Case study applying supervised contrastive learning (Khosla et al., 2020) to facial expression recognition, comparing SupCon loss against a standard cross-entropy baseline in PyTorch with face detection and data augmentation.}
        \resumeItem{Analyzed learned embeddings with t-SNE to diagnose class separation issues and delivered a reproducible codebase to the lab.}
      \resumeItemListEnd

    \resumeSubheading
      {Teaching Assistant, Machine Learning}{Jan 2023 -- Apr 2023}
      {Sathyabama Institute of Science \& Technology}{Chennai, India}
      \resumeItemListStart
        \resumeItem{Graded assignments and homework; prepared lecture slides and supplementary notes for the undergraduate machine learning course.}
        \resumeItem{Led review sessions on mathematics for machine learning, covering matrix algebra and multivariable calculus.}
      \resumeItemListEnd

    \resumeSubheading
      {Summer Intern}{Jun 2022 -- Jul 2022}
      {The University of Texas at Dallas}{Dallas, TX}
      \resumeItemListStart
        \resumeItem{Intensive summer program in machine learning and deep reinforcement learning: lectures, labs, and hands-on model building from first principles.}
      \resumeItemListEnd

  \resumeSubHeadingListEnd

%-----------SKILLS-----------
\section{Technical Skills}
 \begin{itemize}[leftmargin=0.15in, label={}]
    \small{\item{
     \textbf{Languages}{: Python, MATLAB, R, SQL, Bash, LaTeX} \\
     \textbf{Scientific ML}{: PyTorch, NumPy/SciPy, Operator Learning, Physics-Informed Neural Networks (PINNs)} \\
     \textbf{Tools}{: Git, Docker, Linux, HPC job scripting, VS Code, Cursor}
    }}
 \end{itemize}

%-----------READING GROUPS-----------
\section{Reading Groups}
  \resumeSubHeadingListStart
    \resumeSubheading
      {Scientific Machine Learning Reading Group}{2025 -- Present}
      {University of Delaware \, (Advisor: Dr.\ Ke Chen)}{Newark, DE}
      \resumeItemListStart
        \resumeItem{Weekly seminar with rotating student presentations of research literature; topics include Computational Optimal Transport (Peyr\'e \& Cuturi), Equivariant Neural Networks, and One-Shot Learning.}
      \resumeItemListEnd
  \resumeSubHeadingListEnd

\end{document}
